\subsection{ORDERED SUBSETS. PRODUCT OF ORDERED SETS}

Let E be a set ordered by an ordering \(\Gamma\) , with graph G. For each subset A of E, \(G \cap (A \times A)\) is an ordering on A; the corresponding order relation is equivalent to `` \(x \leqslant y\) and \(x \in A\) and \(y \in A\) '', and we shall denote it simply by \(x \leqslant y\) (by abuse of language). The ordering and the order relation thus defined on A are said to be induced by the ordering and order relation given on E; and the ordering and order relation on E are said to be extensions of the ordering and order relation which they induce on A. Whenever we consider A as an ordered set we have in mind the ordering induced on A by that on E, unless the contrary is expressly stated.

Examples. The relations induced by the inclusion relation \(X \subset Y\) on various sets of subsets are of considerable importance. Here are some examples:

\begin{enumerate}
\def\labelenumi{(\arabic{enumi})}
\item
  Let E, F be two sets, and let \(\Phi(E, F)\) be the set of all mappings of subsets of E into F. For each \(f \in \Phi(E, F)\) let \(G_f\) be the graph of \(f\) , which is a subset of \(E \times F\) . If we endow \(\Phi(E, F)\) with the order relation `` \(g\) extends \(f\) '' between \(f\) and \(g\) (no. 1, Example 3), then \(f \to G_f\) is an isomorphism of the ordered set \(\Phi(E, F)\) onto a subset of \(\mathfrak{P}(E \times F)\) , ordered by inclusion.
\item
  For each partition \(\varpi\) of a set E, let \(\tilde{\varpi}\) be the graph of the equivalence relation defined by \(\varpi\) on E. The mapping \(\varpi \to \tilde{\varpi}\) is an isomorphism of the set \(\mathfrak{L}\) of partitions of E, ordered by the relation `` \(\varpi\) is finer than \(\varpi'\) '' between \(\varpi\) and \(\varpi'\) (no. 1, Example 4) onto a subset of \(\mathfrak{P}(E \times E)\) , ordered by inclusion.
\item
  Let E be a set and let \(\Omega \subset \mathfrak{P}(\mathrm{E} \times \mathrm{E})\) be the set of graphs of preorderings on E (no. 2) (or, by abuse of language, the set of all preorderings on E). The order relation \(s \subset t\) between \(s\) and \(t\) , induced on \(\Omega\) by the inclusion relation on \(\mathfrak{P}(\mathrm{E} \times \mathrm{E})\) , is expressed by saying that ``the preordering \(s\) is finer than \(t\) (or that''t is coarser than \(s\) ``). Let \(x(s)y\) and \(x(t)y\) respectively denote the preorder relations \((x,y) \in s\) and \((x,y) \in t\) on E; then to say that \(s\) is finer than \(t\) is equivalent to saying that the relation \(x(s)y\) implies \(x(t)y\) .
\end{enumerate}

Let \((\mathbf{E}_t)_{t \in I}\) be a family of sets, and for each index \(t \in I\) let \(\Gamma_t\) be an ordering on \(\mathbf{E}_t\) ; let \(\mathbf{G}_t \subset \mathbf{E}_t \times \mathbf{E}_t\) be its graph, and let \(x_t \leqslant y_t\) denote the order relation \((x_t, y_t) \in \mathbf{G}_t\) on \(\mathbf{E}_t\) . On the product set \(\mathbf{F} = \prod_{i \in I} \mathbf{E}_t\) , the relation

\[
(\forall \iota) ((\iota \in I) \Longrightarrow (x _ {\iota} \leqslant y _ {\iota}))
\]

is an order relation between \(x = (x_{t})\) and \(y = (y_{t})\) , as is immediately verified. The ordering and the order relation so defined on F are called the product of the orderings \(\Gamma_{t}\) and the product of the order relations \(x_{t} \leqslant y_{t}\) ; this relation is written \(x \leqslant y\) , and the set F, ordered by the product of the orderings \(\Gamma_{t}\) , is called the product of the ordered sets \(E_{t}\) .

It is immediately verified that the graph of the product ordering on \(\mathbf{F}\) is the image of the product set \(\prod_{i\in I}G_i\) under the canonical mapping of \(\prod_{i\in I}(E_i\times E_i)\) onto \(\mathbf{F}\times \mathbf{F}\) (Chapter II, § 5, no. 5).

An important example of a product of ordered sets is the set \(F^{E}\) of graphs of mappings of a set E into an ordered set F. There is a canonical bijection of \(F^{E}\) onto the set \(\mathcal{F}(E, F)\) of mappings of E into F, and this mapping is an isomorphism of the ordered set \(F^{E}\) onto \(\mathcal{F}(E, F)\) endowed with the ordering defined by the relation ``for all \(x \in E\) , \(f(x) \leqslant g(x)\) '' between two mappings f, g of E into F. This relation is written \(f \leqslant g\) .

It should be observed that in the ordered set \(\mathcal{F}(\mathrm{E},\mathrm{F})\) , the relation \(f < g\) means

\pandocbounded{\includegraphics[keepaspectratio]{images/part002_2a386891d62f089ad1c0e080d8072868eea932699f46973e31691c72a41ba0df.jpg}}

``for all \(x \in E\) , \(f(x) \leqslant g(x)\) , and there

and not

\[
f (y) <   g (y) ^ {\prime \prime}
\]

\[
" \text { for   all } x \in \mathrm{E}, f (x) <   g (x)".
\]

In order to avoid confusion it is better not to use the notation \(f < g\) in this situation.

The definitions of this subsection apply without change to preordered sets when ``ordering'' is replaced by ``preordering'' throughout.

